\section{Zadanie 3}
Udowodnij, że jeżeli jeden ze stanów w danej klasie relacji skomunikowania
jest stanem powracającym, to wszystkie stany w tej klasie są powracające.
\subsection*{Rozwiązanie}
Tak naprawdę cała intuicja jest taka że, skoro rozpatrujemy daną klasę skomunikowania to zachodzi: $\forall_{i,j} i\leftarrow j \text{ oraz } j\leftarrow i$. Załóżmy że wiemy, że stanem powracającym jest $i$. Teraz skoro to jest klasa skomunikowania to jak weźmiemy jakiś stan $x$ to zachodzi $i\leftarrow x$, ale ze względu na klasę skomunikowania zachodzi również $x\leftarrow i$. Zatem zauważmy że skoro istnieje ścieżka z $i$ do $i$ przez $x$ oraz $i$ jest powracający to $x$ też jest powracający.\\
Bardziej formalnie:\\

Jeśli $i$ i $j$ są w tej samej klasie komunikacji i stan $i$ jest powracający,
to stan $j$ również jest powracający.
\par
Ponieważ $i \leftrightarrow j$, istnieją liczby całkowite $n_0,m_0 \ge 1$
takie, że
\[
p_{j i}(n_0) > 0
\qquad\text{oraz}\qquad
p_{i j}(m_0) > 0 .
\]

Założenie, że stan $i$ jest powracający oznacza, że
\[
\sum_{r=0}^{\infty} p_{ii}(r) = +\infty.
\]

Dla każdego $r \ge 0$ zachodzi
\[
p_{j j}(n_0 + r + m_0)
\;\ge\;
p_{j i}(n_0) \, p_{i i}(r) \, p_{i j}(m_0).
\]

Sumując obie strony po wszystkich $r \ge 0$ otrzymujemy
\[
\sum_{t = 0}^{\infty} p_{j j}(t)
\;\ge\;
p_{j i}(n_0) \, p_{i j}(m_0)
\sum_{r = 0}^{\infty} p_{i i}(r).
\]

Ponieważ $p_{ji}(n_0) > 0$, $p_{ij}(m_0) > 0$ oraz
$\sum_{r=0}^{\infty} p_{ii}(r) = +\infty$, prawa strona sumy jest równa
$+\infty$. Zatem
\[
\sum_{t=0}^{\infty} p_{jj}(t) = +\infty,
\]
co oznacza, że stan $j$ jest powracający.

Wynika stąd, że każdy stan w klasie komunikacji, która zawiera stan powracający,
również jest powracający.
